% !PRINTED 69
\weberpage{69}

\begin{center}
\textbf{\S\,20. \gesperrt{Die $\vartheta$-Funktion.}}
\end{center}

Die Funktion $\vartheta(v)$ ist eine \gesperrt{ganze} Funktion von $v$, die den
Bedingungen gen\"ugt:

\weq{1}{\begin{aligned}
\vartheta(v + 1) &= e^{-\pi i g_1}\vartheta(v),\\
\vartheta(v + \omega) &= e^{-\pi i(2v+\omega)} e^{-\pi i g_2}\vartheta(v),
\end{aligned}}

\gesperrt{und ist dadurch bis auf einen von $v$ unabh\"angigen Faktor
bestimmt.} Sie ist also unter den $t$-Funktionen als Spezialfall
enthalten, w\"ahrend andererseits die allgemeine $t$-Funktion durch
die $\vartheta$-Funktion ausgedr\"uckt werden kann in der Weise:
\[
t(u) = C\, e^{-\pi i(\lambda u^2 + \mu u)}\, \vartheta\!\left(\frac{u}{\omega_1},\ \frac{\omega_2}{\omega_1}\right),
\]
worin $C, \lambda, \mu$ beliebige konstante oder von $\omega_1, \omega_2$ abh\"angige
Gr\"o\ss en sind.

Die Funktion $\vartheta$ ist noch von den in der Charakteristik vorkommenden
 Zahlen $g_1, g_2$ abh\"angig, und wenn eine Bezeichnung
dieser Abh\"angigkeit erforderlich ist, so soll f\"ur $\vartheta(v, \omega)$:
\[
\vartheta_{g_1,\, g_2}(v, \omega)
\]
gesetzt werden, es k\"onnen dabei $g_1, g_2$ beliebige Zahlen sein.
Nach dem fr\"uheren gen\"ugt es, wenn wir sie reell und zwischen 0
und 2 annehmen.

Entsprechend den $T$-Funktionen h\"oherer Ordnung werden
wir auch $\Theta$-Funktionen $m$ter Ordnung einf\"uhren und verstehen
darunter eine ganze Funktion von $v$, die den Bedingungen gen\"ugt:

\weq{2}{\begin{aligned}
\Theta(v + 1) &= e^{-\pi i g_1}\Theta(v)\\
\Theta(v + \omega) &= e^{-m\pi i(2v+\omega)} e^{-\pi i g_2}\Theta(v).
\end{aligned}}

Auch hierbei kann die Charakteristik $g_1, g_2$ in die Bezeichnung
 mit aufgenommen werden:
\[
\Theta_{g_1,\, g_2}(v, \omega).
\]

Diese Funktionen sind als spezielle F\"alle unter den $T$-Funktionen
 enthalten, und man kann allgemeine $T$-Funktionen in der
Weise bilden:

\weq{3}{C\, e^{-\pi i(\lambda u^2 + \mu u)}\, \Theta\!\left(\frac{u}{\omega_1},\ \frac{\omega_2}{\omega_1}\right).}

Es ergibt sich also aus \S\,18, 3. der \gesperrt{Fundamentalsatz} f\"ur
die $\Theta$-Funktionen, den wir folgenderma\ss en aussprechen:
%CONT>>
